\documentclass[11pt]{article}
\usepackage[margin=1in]{geometry}
\usepackage[hidelinks]{hyperref}

\begin{document}

\begin{titlepage}
\centering
{\Large George Washington University Data Science Program\par}
\vspace{1cm}
{\large Capstone Report --- Fall 2026\par}
\vspace{1.5cm}
{\huge\bfseries Exploring Applications for the Montandon Global Crisis Data Bank\par}
\vspace{2cm}
{\large Jeongmin An, Jehan Bugli, and Aidan Carlisle\par}
\vspace{1cm}
{\large Advisor: Amir Jafari\par}
\vfill
{\large Draft --- September 2026\par}
\end{titlepage}

\tableofcontents
\newpage

\section{Introduction}

This report documents a project exploring applications of the Montandon Global Crisis
Data Bank. The project combines disaster and response records with external news,
geospatial context, text embeddings, and network analysis.

\section{Problem Statement}

Montandon contains rich crisis and response information, but connecting related records
and exploring their relationships can be difficult. This project investigates whether a
graph-based representation and a simple query interface can make that information easier
to explore. Results and evaluation criteria will be added as the project is scoped.

\section{Related Work}

Crisis risk reduction depends on understanding disaster risk, preparedness, and recovery
\cite{undrr2015sendai}. Research on emergency social-media messages
includes filtering, classification, and summarization \cite{imran2015processing}; this
project instead starts with the curated Montandon data bank and uses external articles as
additional context. STAC provides a common structure for spatiotemporal metadata
\cite{ogc2025stac}, which is relevant because Montandon uses a modified STAC approach.

\section{Solution and Methodology}

This section will describe data access, feature generation, graph construction, and the query setup.

\section{Results and Discussion}

This section will report the project’s results, figures, and limitations.

\section{Conclusion}

This section will summarize findings and recommended next steps.

\bibliographystyle{IEEEtran}
\bibliography{references}

\end{document}
